\documentclass[12pt]{scrartcl}

\input{../styles/Packages}
\input{../styles/FormatAndHeader}

\usepackage{xcolor} % Für die Farben im Code

% Definiere das Aussehen des Codes
\lstset{
    language=SQL,
    basicstyle=\ttfamily\small,
    keywordstyle=\color{blue}\bfseries,
    stringstyle=\color{red},
    commentstyle=\color{gray}\itshape,
    numbers=left,           % Zeilennummern links
    numberstyle=\tiny\color{gray},
    stepnumber=1,
    showstringspaces=false, % Keine komischen Unterstriche in Strings
    tabsize=4,
    breaklines=true,        % Automatischer Zeilenumbruch
    breakatwhitespace=false,
    frame=single,           % Ein Rahmen um den Code
    extendedchars=true,
    literate={ö}{{\"o}}1 {ä}{{\"a}}1 {ü}{{\"u}}1 {ß}{{\ss}}1 {Ö}{{\"O}}1 {Ä}{{\"A}}1 {Ü}{{\"U}}1
}

% \stepcounter{countername}: Increases counter
\setcounter{sheetnr}{3} % Number of sheet


\setcounter{exnum}{1} % Exercise number

\begin{document}
% Nutze den \exercise{Aufgabenname} Befehl, um eine neue Aufgabe zu beginnen.
% Möchtest du eine Aufgabe in der Nummerierung überspringen, schreibe vor der Aufgabe: \stepcounter{exnum}
% Möchtest du die Nummer einer Aufgabe auf eine beliebige Zahl x setzen, schreibe vor der Aufgabe: \setcounter{exnum}{x}

    %TODO
    \exercise{Abgabe 3.1: SQL-Abfragen}


    \newpage

    \exercise{Abgabe 3.2: SQL-Anweisungen}

    \subsection*{a) Neue Abteilung in München eröffnen}
% Mit dem INSERT INTO Befehl fügen wir einen neuen Datensatz in die Tabelle Abteilung ein.
    \begin{lstlisting}[language=SQL]
INSERT INTO Abteilung (AbteilungsNr, Name, Ort)
VALUES (11, 'A11', 'München');
    \end{lstlisting}

    \subsection*{b) Alle Abteilungen in Stuttgart löschen}
% DELETE Anweisung mit einfacher Filterung auf den Ort.
    \begin{lstlisting}[language=SQL]
DELETE FROM Abteilung
WHERE Ort = 'Stuttgart';
    \end{lstlisting}

    \subsection*{c) Mitarbeiter aus Köln in Abteilung 11 versetzen}
% Standard-SQL sieht keinen direkten JOIN in der UPDATE-Klausel vor.
% Daher ermitteln wir die Kölner Abteilungen über eine unkorrelierte Unterabfrage.
    \begin{lstlisting}[language=SQL]
UPDATE Angestellter
SET AbteilungsNr = 11
WHERE AbteilungsNr IN (
    SELECT AbteilungsNr
    FROM Abteilung
    WHERE Ort = 'Köln'
);
    \end{lstlisting}

    \subsection*{d) Sicht für Hamburger Projekte mit mind. 10 Mitarbeitern}
% Joins laut Vorgabe über die WHERE-Klausel. Um zu zählen, müssen wir gruppieren.
% In der SELECT-Klausel dürfen nur Attribute stehen, die auch gruppiert wurden
% (oder aggregiert werden).
    \begin{lstlisting}[language=SQL]
CREATE VIEW HamburgerProjekte AS
SELECT Projekt.ProjektNr, Projekt.Projektname, Projekt.Projektort
FROM Projekt, Projektmitarbeit
WHERE Projekt.ProjektNr = Projektmitarbeit.ProjektNr
  AND Projekt.Projektort = 'Hamburg'
GROUP BY Projekt.ProjektNr, Projekt.Projektname, Projekt.Projektort
HAVING COUNT(Projektmitarbeit.PersonalNr) >= 10;
    \end{lstlisting}

    \subsection*{e) Inhalt der Tabelle Angestellter löschen}
% Ein DELETE ohne WHERE-Klausel löscht alle Datensätze, lässt die
% Tabellenstruktur aber unberührt.
    \begin{lstlisting}[language=SQL]
DELETE FROM Angestellter;
    \end{lstlisting}

    \newpage

    \subsection*{f) Tabellen Angestellter und Abteilung anlegen}
% Die Datentypen, Längen und NOT NULL Constraints (nullable = false) wurden exakt
% aus der ODPS-Spezifikation übernommen.
% Die Tabelle 'Abteilung' wird zuerst angelegt, da 'Angestellter' darauf referenziert.
\begin{lstlisting}[language=SQL]
CREATE TABLE Abteilung (
    AbteilungsNr INTEGER NOT NULL,
    Name VARCHAR(25) NOT NULL,
    Ort VARCHAR(25),
    PRIMARY KEY (AbteilungsNr)
);

CREATE TABLE Angestellter (
    PersonalNr INTEGER NOT NULL,
    Name VARCHAR(25),
    Vorname VARCHAR(25),
    Beruf VARCHAR(25),
    Gehalt DECIMAL(8),
    AbteilungsNr INTEGER NOT NULL,
    Manager INTEGER,
    Geburtstag DATE,
    Wohnort VARCHAR(25),
    PRIMARY KEY (PersonalNr),
    FOREIGN KEY (AbteilungsNr) REFERENCES Abteilung(AbteilungsNr)
        ON DELETE CASCADE
        ON UPDATE RESTRICT,
    FOREIGN KEY (Manager) REFERENCES Angestellter(PersonalNr)
);
\end{lstlisting}


\end{document}
